\renewcommand{\chapterprefix}{Appendix}\renewcommand{\thechapter}{S}
\chapter{Solutions to the NOW YOU exercises}
\chaptercontents{A & Chapter 0 \\ B & Chapter 1 \\ C & Chapter 2 \\ D & Chapter 3 \\ E & Chapter 4 \\ F & Chapter 5}{Every book exercise in Chapters 0--5 that is not worked in the main text is solved here, in the order of the book. Compare line by line: did you introduce every variable, unfold every definition, and name the strategy?}

\hsection{Chapter 0}
\sol{0.2} (a) $2+2=4$. (b) $2+2=5$. (c) ``There are infinitely many pairs of primes that differ by $2$'' (the twin prime conjecture; nobody knows its truth value). Any proposition whose truth value is unknown is fine, e.g.\ $\pi^\pi\in\Q$ (Exercise 0.21(f)).

\sol{0.47} Repeated division by $17$: $408735787=24043281\cdot17+10$, $24043281=1414310\cdot17+11$, $1414310=83194\cdot17+12$, $83194=4893\cdot17+13$, $4893=287\cdot17+14$, $287=16\cdot17+15$, $16=0\cdot17+16$. Reading the remainders upwards, with $A=10,\dots,G=16$: $408\,735\,787=GFEDCBA_{(17)}$. Likewise $1442151747=40059770\cdot36+27$, $40059770=1112771\cdot36+14$, $1112771=30910\cdot36+11$, $30910=858\cdot36+22$, $858=23\cdot36+30$, $23=0\cdot36+23$, and $27=R$, $14=E$, $11=B$, $22=M$, $30=U$, $23=N$: $1\,442\,151\,747=NUMBER_{(36)}$.

\sol{0.55} $s(w)=\sqrt2\,w^4-81\sqrt2$ has coefficients $\sqrt2,-81\sqrt2\notin\Q$ (so it is not a polynomial over $\Q$), is quartic, and $s(3)=\sqrt2(81-81)=0\in\Q$.

\sol{0.64} Since $i^2=-1$, $i^3=-i$ and $i^4=1$. For $n\in\Z$ write $n=4q+r$ with $r\in\{0,1,2,3\}$ (Theorem 0.42); then $i^n=(i^4)^qi^r=i^r$. So $i^n=1$ if $r=0$; $i^n=0+1\cdot i$ if $r=1$; $i^n=-1$ if $r=2$; $i^n=0+(-1)i$ if $r=3$. (For negative $n$ this uses $i^{-1}=-i$, which is consistent with the formula.)

\sol{0.65} Let $p(x)=ax^2+bx+c$ with $a\ne0$ and $\Delta<0$. By Definition 0.62, $\sqrt\Delta=\sqrt{|\Delta|}\,i$, so $(\sqrt\Delta)^2=|\Delta|i^2=-|\Delta|=\Delta$. The computation in the proof of Theorem 0.60 used only $(\sqrt\Delta)^2=\Delta$ and the completed square $p(x)=a(x+\frac b{2a})^2-\frac\Delta{4a}$, so $p(\alpha)=p(\beta)=0$ holds verbatim. For ``no other roots'': if $p(\gamma)=0$ then $(\gamma+\frac b{2a})^2=\frac{\Delta}{4a^2}=\big(\frac{\sqrt\Delta}{2a}\big)^2$, so $\big(\gamma+\frac b{2a}-\frac{\sqrt\Delta}{2a}\big)\big(\gamma+\frac b{2a}+\frac{\sqrt\Delta}{2a}\big)=0$, and a product of complex numbers is $0$ only if a factor is $0$; hence $\gamma=\alpha$ or $\gamma=\beta$. The cases $\Delta\ge0$ are Theorem 0.60.

\sol{0.66} Let $\alpha=a+bi$ and put $p(x)=x^2-2ax+(a^2+b^2)$, a quadratic over $\R$. Then $p(\alpha)=(a+bi)^2-2a(a+bi)+a^2+b^2=a^2+2abi-b^2-2a^2-2abi+a^2+b^2=0$. The other root is $a-bi$: $p(a-bi)=a^2-2abi-b^2-2a^2+2abi+a^2+b^2=0$. (Indeed $p(x)=(x-(a+bi))(x-(a-bi))$.)

\sol{0.67} \emph{Quadratic formula over $\C$.} Let $p(x)=ax^2+bx+c$ with $a,b,c\in\C$, $a\ne0$, and let $\delta\in\C$ be any number with $\delta^2=b^2-4ac$ (every complex number has a square root; we take this for granted). Then the roots of $p(x)$ are precisely $\frac{-b+\delta}{2a}$ and $\frac{-b-\delta}{2a}$. \emph{Proof:} completing the square, $p(x)=a\big(x+\frac b{2a}\big)^2-\frac{\delta^2}{4a}$; substituting the two values gives $a\cdot\frac{\delta^2}{4a^2}-\frac{\delta^2}{4a}=0$; and if $p(\gamma)=0$ then $(\gamma+\frac b{2a})^2=(\frac\delta{2a})^2$, so $\gamma+\frac b{2a}=\pm\frac\delta{2a}$ by factorising the difference of squares.

\sol{0.68} \emph{Subtraction} $a-b$: translate $a$ by $b$ units to the \emph{left} (it is the addition of $-b$): put the $0$ of the lower line under $a$ and read $a-b$ above the point $-b$ of the lower line. \emph{Division} $a/b$ ($b\ne0$): it is the inverse of multiplication. Align the zeros, then stretch (and reflect if $b<0$) the lower line until its point $b$ lies beneath the point $a$ of the upper line; then $a/b$ is the point of the upper line above the point $1$ of the lower line (because the scaling sends $1\mapsto a/b$ and $b\mapsto a$).

\sol{0.E.1} With $A,\dots,Z\mapsto0,\dots,25$; $a,\dots,z\mapsto26,\dots,51$; $0,\dots,9\mapsto52,\dots,61$; $\text{-}\mapsto62$; $\_\mapsto63$, the digits of dQw4w9WgXcQ are
\[
29,\ 16,\ 48,\ 56,\ 48,\ 61,\ 22,\ 32,\ 23,\ 28,\ 16,\qquad\text{and}\qquad \sum_kd_k64^k=33\,736\,714\,463\,647\,725\,328 .
\]
Repeated division of $7\,159\,047\,702\,620\,056\,984$ by $64$ gives the remainders $24,28,2,29,45,18,3,44,26,13,6$; reading them last-first, $6,13,26,44,3,18,45,29,2,28,24$ are the characters G, N, a, C, \_, S, M, v, q, 2, Y, so the base-64 expansion is the string GNaC\_SMvq2Y.

\sol{0.E.2} $(a+b\sqrt2)(c+d\sqrt2)=(ac+2bd)+(ad+bc)\sqrt2$. If $ad+bc=0$ the product is $ac+2bd\in\Z$. If $ad+bc\ne0$ and the product were an integer $m$, then $\sqrt2=\frac{m-(ac+2bd)}{ad+bc}\in\Q$, contradicting Assumption 0.49. So the product is an integer exactly when $ad+bc=0$.

\sol{0.E.6--0.E.11} $a\odot b=a+b$: $\N,\Z,\Q,\R$ all closed (Assumptions 0.7, 0.9, Theorem 0.16, Assumption 0.5). \ $a-b$: $\N$ not closed ($1-2=-1\notin\N$); $\Z,\Q,\R$ closed. \ $(a-b)(a+b)$: $\N$ not closed ($0\odot1=-1$); $\Z,\Q,\R$ closed (products and sums). \ $(a-1)(b-1)+2(a+b)=ab+a+b+1$: all four closed (sums and products only). \ $\frac a{b^2+1}$: $\N$ and $\Z$ not closed ($1\odot1=\frac12$); $\Q$ and $\R$ closed, as $b^2+1\ne0$. \ $\frac a{\sqrt{b^2+1}}$: $\N,\Z,\Q$ not closed, since $1\odot1=\frac1{\sqrt2}$ is irrational (if $\frac1{\sqrt2}=r\in\Q$ then $r\ne0$ and $\sqrt2=\frac1r\in\Q$); $\R$ closed, as $\sqrt{b^2+1}$ is a non-zero real.

\sol{0.E.13} Let $x\in\Q$. The polynomial $p(t)=t-x$ is non-zero, has coefficients $1,-x\in\Q$, and $p(x)=0$. So $x$ is algebraic.

\sol{0.E.14} $p(t)=t^2-2$ is non-zero over $\Q$ and $p(\sqrt2)=2-2=0$.

\sol{0.E.15} Let $x=\sqrt2+\sqrt3$. Then $x^2=5+2\sqrt6$, so $(x^2-5)^2=24$, i.e.\ $x^4-10x^2+1=0$. So $x$ is a root of the non-zero polynomial $t^4-10t^2+1$ over $\Q$.

\sol{0.E.16} By Exercise 0.66, $x+yi$ is a root of $t^2-2xt+(x^2+y^2)$, whose coefficients $1,-2x,x^2+y^2$ are rational when $x,y\in\Q$ (Theorem 0.16), and which is non-zero. So $x+yi$ is algebraic.

\sol{0.E.19} True: $0=0\cdot n$ for every $n\in\Z$ (Exercise 0.38).

\sol{0.E.20} True: $n^2=n\cdot n$ with $n\in\Z$, so $n$ divides $n^2$.

\sol{0.E.21} True: if $x\in\Q$ then $x^2=x\cdot x\in\Q$ by Theorem 0.16.

\sol{0.E.23} False: $a=0$, $b=5$: $0=0\cdot5+0$, remainder $0$, which is not less than $a=0$. (Also any negative $a$: $-3=(-1)\cdot5+2$ and $2>-3$.)

\sol{0.E.24} False: $x^2$ has discriminant $0$ and its only root is $0$ (if $\gamma^2=0$ then $\gamma=0$), so it does not have two distinct complex roots.

\sol{0.E.25} Always. Since $n<b_1<b_2$, condition (ii) allows $n$ itself as a single digit in both bases, and $n=n\cdot b^0$ with $n\ne0$, so the base-$b_1$ expansion and the base-$b_2$ expansion of $n$ are both the one-digit string $n$.

\sol{0.E.26} Never. Suppose the two expansions are the same string $d_r\dots d_0$. Since $n>b_1$, the base-$b_1$ expansion has at least two digits, so $r\ge1$ and $d_r\ne0$. Then $n=\sum_{i\le r}d_ib_1^i$ and also $n=\sum_{i\le r}d_ib_2^i$. But $b_1<b_2$ gives $d_rb_1^r<d_rb_2^r$ and $d_ib_1^i\le d_ib_2^i$ for the other terms, so $\sum d_ib_1^i<\sum d_ib_2^i$, i.e.\ $n<n$: contradiction.

\sol{0.E.31} Sometimes: $x=1$ has $x^2=1\in\Q$ and $x\in\Q$; $x=\sqrt2$ has $x^2=2\in\Q$ but $x\notin\Q$.

\sol{0.E.32} Always. From $x^2+1\in\Q$ we get $x^2\in\Q$, hence $x^4=(x^2)^2\in\Q$ (Theorem 0.16); from $x^5+1\in\Q$ we get $x^5\in\Q$. If $x=0$ then $x\in\Q$. If $x\ne0$ then $x^4\ne0$ and $x=\frac{x^5}{x^4}\in\Q$ by Theorem 0.16.

\sol{0.E.33} Always. Since $u+vi$ with $v\ne0$ is not real, $p(x)$ has a non-real root, so $\Delta<0$ (by Exercise 0.61 and Theorem 0.60, $\Delta\ge0$ gives only real roots). By Exercise 0.65 the roots are $\frac{-b\pm\sqrt{|\Delta|}\,i}{2a}$, which are of the form $u'\pm v'i$ with $u'=\frac{-b}{2a}$, $v'=\frac{\sqrt{|\Delta|}}{2a}$. One of them is $u+vi$, so $u=u'$ and $v=\pm v'$, and the other root is $u-vi$.

\hsection{Chapter 1}
\sol{1.1.6} $p$: ``John is a mathematician''; $q$: ``John lives in Pittsburgh''; the statement is $p\wedge q$.

\sol{1.1.13} Assume $(p\wedge q)\wedge r$. By Strategy 1.1.10 we may assume $p\wedge q$ (and $r$). By Strategy 1.1.10 again, from $p\wedge q$ we may assume $q$. Hence $q$ is true.

\sol{1.1.14} Proof tree: from $p$ and $q$ conclude $p\wedge q$ ($\wedge$I); from $r$ and $s$ conclude $r\wedge s$ ($\wedge$I); from $p\wedge q$ and $r\wedge s$ conclude $(p\wedge q)\wedge(r\wedge s)$ ($\wedge$I). With $p$: ``$2$ is even'', $q$: ``$2$ is prime'', $r$: ``$3$ is odd'', $s$: ``$3$ is prime'', a proof has four parts (prove $2$ even, prove $2$ prime, prove $3$ odd, prove $3$ prime), combined two at a time.

\sol{1.1.23} Since $n=8$ or $n=9$, we split into cases. Case $n=8$: $8=2\cdot4$, so $n$ is even. Case $n=9$: $9=3^2$, so $n$ is a perfect square. In both cases, $n$ is even or a perfect square.

\sol{1.1.26} The discriminant is $\Delta=a^2-4b\ne0$. If $\Delta>0$, Theorem 0.60 gives $\alpha=\frac{-a+\sqrt\Delta}2$, $\beta=\frac{-a-\sqrt\Delta}2$ (in some order), so $\alpha-\beta=\pm\sqrt\Delta$; take $c=\pm\sqrt\Delta\in\R$. If $\Delta<0$, Exercise 0.65 gives the roots $\frac{-a\pm\sqrt{|\Delta|}\,i}2$, so $\alpha-\beta=\pm\sqrt{|\Delta|}\,i=ci$ with $c=\pm\sqrt{|\Delta|}\in\R$. Either way such a $c$ exists.

\sol{1.1.31} Suppose $\alpha$ is a root of $p(x)$ and $a\in\C$. Let $q(x)=(x-a)p(x)$. Then $q(\alpha)=(\alpha-a)p(\alpha)=(\alpha-a)\cdot0=0$, so $\alpha$ is a root of $q(x)$.

\sol{1.1.64} In the proof of Proposition 1.1.63, the law of excluded middle is used to split into the cases ``$a$ is even'' and ``$a$ is odd'' (odd means not even). Proof by contradiction is used inside the second case: assuming $b$ odd we derived that $ab$ is odd, which contradicts the hypothesis that $ab$ is even; the contradiction is ``$ab$ is even and $ab$ is not even''. \emph{Version without LEM:} suppose $ab$ is even, and suppose towards a contradiction that it is not the case that $a$ or $b$ is even; then (de Morgan) $a$ is odd and $b$ is odd, so $a=2k+1$, $b=2\ell+1$ and $ab=2(2k\ell+k+\ell)+1$ is odd, contradicting $ab$ even. \emph{Version with LEM and no contradiction:} by LEM, $a$ is even or odd, and $b$ is even or odd. If $a$ is even or $b$ is even we are done. In the remaining case both are odd, and then $ab=2(2k\ell+k+\ell)+1$ leaves remainder $1$ on division by $2$, so by the uniqueness clause of the division theorem $ab$ is not of the form $2m$; this case is therefore excluded by the hypothesis that $ab$ is even. (Any proof about ``odd'' numbers must use negation somewhere, because odd is defined as not even.)

\sol{1.2.8} (a) ``$b$ divides $a$.'' (b) ``$x$ is rational.'' (c) ``The only natural numbers that divide $n$ are $1$ and $n$'' (true for $n=1$ and for primes). (d) ``There is no largest positive real number'' (every positive real is exceeded by a positive real).

\sol{1.2.14} Let $x\in\R$. Every object is equal to itself (the basic assumption mentioned in Definition 0.3), so $x=x$. Since $x$ was arbitrary, $\forall x\in\R,\ x=x$.

\sol{1.2.16} Let $n\in\Z$. Then $n=\frac n1$ with $n,1\in\Z$ and $1\ne0$, so $n\in\Q$ by Definition 0.15.

\sol{1.2.17} Let $p(x)=ax+b$ be a linear polynomial over $\Q$, so $a,b\in\Q$ and $a\ne0$. Define $\gamma=-\frac ba$; then $\gamma\in\Q$ by Theorem 0.16, and $p(\gamma)=-b+b=0$. So $p(x)$ has a rational root.

\sol{1.2.27} Define $x=\sqrt2$. Then $x$ is irrational (Assumption 0.49) and $x^2=2=\frac21$ is rational.

\sol{1.2.34} The existence part gives at least one $x$ with $p(x)$; the uniqueness part says any two elements satisfying $p$ are equal, so there cannot be two different ones; together: exactly one. Alternative definitions: Proposition 1.2.38's form $\exists x\in X,\ (p(x)\wedge\forall y\in X,\ (p(y)\Rightarrow y=x))$; or ``$\exists x,\ p(x)$ and $\neg\exists a,\exists b,\ (a\ne b\wedge p(a)\wedge p(b))$''.

\sol{1.2.37 (c)} $\exists!n\in\N,\ (n\text{ has exactly one positive divisor})$. \emph{Existence:} $n=1$: its positive divisors are the positive $d$ with $1=qd$, and then $d=1$ (if $d\ge2$ then $qd\ne1$ for $q\in\Z$); so exactly one. \emph{Uniqueness:} let $n\in\N$ have exactly one positive divisor. By Exercise 0.38, $1$ divides $n$, so that divisor is $1$. If $n\ge2$ then $n$ also divides $n$ and $n\ne1$: two positive divisors, contradiction. If $n=0$ then every positive integer divides $0$: more than one, contradiction. Hence $n=1$. So any two such natural numbers are both $1$.

\sol{1.3.9} $\neg p$: by ($\neg$I)/($\neg$E), $\neg p$ is true exactly when $p$ is false. $p\wedge q$: by ($\wedge$I) it is true when both are, and by ($\wedge$E$_1$), ($\wedge$E$_2$) it cannot be true when either is false. $p\vee q$: by ($\vee$I) it is true when at least one is; if both are false, $p\vee q$ cannot be true, since by ($\vee$E) we could then derive $\bot$ (both $p$ and $q$ give contradictions). $p\Rightarrow q$: by ($\Rightarrow$E) it is false when $p$ is true and $q$ false; when $q$ is true, $q$ can be derived from anything, so ($\Rightarrow$I) makes it true; when $p$ is false, assuming $p$ gives a contradiction and then $q$ by explosion, so ($\Rightarrow$I) makes it true.

\sol{1.3.25} \emph{(i) Derivations.} Assume $\neg(p\vee q)$. To prove $\neg p$: assume $p$; then $p\vee q$ ($\vee$I$_1$), contradicting $\neg(p\vee q)$; so $\neg p$. Likewise $\neg q$. Hence $(\neg p)\wedge(\neg q)$. Conversely assume $(\neg p)\wedge(\neg q)$. To prove $\neg(p\vee q)$, assume $p\vee q$ and split: $p$ contradicts $\neg p$; $q$ contradicts $\neg q$. So $\neg(p\vee q)$. \emph{(ii) Truth table:} $\neg(p\vee q)$ is $\checkmark$ only in row $\times\times$, and so is $(\neg p)\wedge(\neg q)$. \emph{(iii)} By part (a) applied to $\neg p,\neg q$: $\neg((\neg p)\wedge(\neg q))\equiv(\neg\neg p)\vee(\neg\neg q)\equiv p\vee q$ (double negation). Negating both sides, $(\neg p)\wedge(\neg q)\equiv\neg\neg((\neg p)\wedge(\neg q))\equiv\neg(p\vee q)$.

\sol{1.3.38} $\neg(p\Leftrightarrow q)\equiv\neg[(p\Rightarrow q)\wedge(q\Rightarrow p)]\equiv\neg(p\Rightarrow q)\vee\neg(q\Rightarrow p)\equiv(p\wedge\neg q)\vee(q\wedge\neg p)$, using Definition 1.1.40, de Morgan (a) and Exercise 1.3.27.

\sol{1.3.44 (b), (d)} (b) Assume $p\Rightarrow(q\Rightarrow r)$; assume $p\Rightarrow q$; assume $p$. Modus ponens gives $q$ (from $p\Rightarrow q$) and $q\Rightarrow r$ (from the first assumption), then $r$. So the formula is derivable from no assumptions, hence a tautology. \emph{Meaning:} an implication can be distributed over an implication. (d) By Theorem 1.3.24(a), $\neg(p\wedge q)\equiv(\neg p)\vee(\neg q)$, so by Theorem 1.3.45(b) the biconditional is a tautology. \emph{Meaning:} to show a conjunction fails, show one conjunct fails.

\sol{1.E.1} $p$: $n$ even, $q$: $n$ prime, $r$: $n=2$. (a) ``If $n$ is an even prime, then $n=2$'': true for all $n$ (an even prime is divisible by $2$, so $2$ is one of its only positive divisors $1$ and $n$, and $2\ne1$). (b) ``If $n$ is a prime other than $2$, then $n$ is odd'': true for all $n$ (same reason). (c) ``$n$ is odd, or $n$ is not prime, or $n\ne2$'': true for every $n$ except $n=2$, for which all three disjuncts fail; so true for some, false for some. (d) ``$n$ is an even prime different from $2$'': false for all $n$.

\sol{1.E.2} (a) $p$: guinea pigs are quiet; $q$: they are hungry; $r$: they are loud: $p\wedge(q\Rightarrow r)$. (b) $p$: $2$ is even; $q$: $2$ is prime: $p\wedge q$. (c) $p$: $\sqrt2$ is irrational; $q$: $\sqrt2$ is an integer: $p\wedge(p\Rightarrow\neg q)$ (hence $\neg q$).

\sol{1.E.3} ``$A$ is necessary for $B$'' means $B\Rightarrow A$; ``$A$ is sufficient for $B$'' means $A\Rightarrow B$. (a) $p\Rightarrow\neg q$: true. (b) $\neg q\Rightarrow p$: false. (c) $\neg q\Rightarrow p$: false. (d) $p\Rightarrow\neg q$: true.

\sol{1.E.4} To prove $(p\wedge q)\Rightarrow\neg r$: assume $p\wedge q$, so assume $p$ and $q$; then prove $\neg r$ by assuming $r$ and deriving a contradiction.

\sol{1.E.5} Assume $p\vee q$; assume $r$; prove $s$ by cases: derive $s$ from $p$ (and $r$), then derive $s$ from $q$ (and $r$).

\sol{1.E.6} Prove both implications. ($\Rightarrow$) Assume $p\Rightarrow q$; assume $\neg p$; prove $\neg q$ by assuming $q$ and deriving a contradiction. ($\Leftarrow$) Assume $\neg p\Rightarrow\neg q$; assume $p$; prove $q$. \emph{Warning:} this formula is not a tautology (with $p$ true and $q$ false, $p\Rightarrow q$ is false but $\neg p\Rightarrow\neg q$ is true), so the second step cannot succeed in general: $\neg p\Rightarrow\neg q$ is the \emph{inverse}, equivalent to the converse $q\Rightarrow p$, not to $p\Rightarrow q$.

\sol{1.E.7} Prove one of the disjuncts: either prove $p$ and prove $\neg q$ (assume $q$, derive a contradiction), or prove $q$ and prove $\neg p$. (Not a tautology: false when $p,q$ are both true.)

\sol{1.E.8} ``Between any two distinct rational numbers there is an irrational number.'' \emph{Proof.} Let $a,b\in\Q$ and suppose $a<b$. Define $c=a+\frac{b-a}{\sqrt2}$. First, $\frac{b-a}{\sqrt2}=\frac{b-a}2\cdot\sqrt2$ is irrational: if it were rational, then $\sqrt2=\frac{b-a}{\sqrt2}\cdot\frac2{b-a}$ would be rational (Theorem 0.16, $b-a\ne0$), contradicting Assumption 0.49. Hence $c$ is irrational by Proposition 1.1.53. Second, since $\sqrt2>1$ we have $0<\frac1{\sqrt2}<1$, so $0<\frac{b-a}{\sqrt2}<b-a$ and therefore $a<c<b$. So $c$ is the required number.

\sol{1.E.9} $\forall x\in\Z,\ \exists y\in\Z,\ y>x$. \emph{Proof.} Let $x\in\Z$. Define $y=x+1$. Then $y\in\Z$ (closure) and $y>x$.

\sol{1.E.10} $\exists a\in X,\ \exists b\in X,\ \big(a\ne b\wedge p(a)\wedge p(b)\wedge\forall x\in X,\ (p(x)\Rightarrow(x=a\vee x=b))\big)$. \emph{Structure:} exhibit two distinct elements satisfying $p$, then show every element satisfying $p$ is one of them. \emph{Proof for $x^2=1$.} Take $a=1$, $b=-1$: distinct, and $1^2=(-1)^2=1$. Let $x\in\R$ with $x^2=1$. Then $(x-1)(x+1)=0$, so $x-1=0$ or $x+1=0$, i.e.\ $x=1$ or $x=-1$.

\sol{1.E.12} The truth tables of $p\Rightarrow q$ and $q\Rightarrow p$ differ in the rows $(\checkmark,\times)$ and $(\times,\checkmark)$: in the first, $p\Rightarrow q$ is $\times$ and $q\Rightarrow p$ is $\checkmark$. Example: $p$: ``$1=2$'', $q$: ``$1=1$''; $p\Rightarrow q$ is true (false hypothesis) and $q\Rightarrow p$ is false.

\sol{1.E.13} The column is $\checkmark$ exactly when $p,q$ differ: $(p\wedge\neg q)\vee(q\wedge\neg p)$ (equivalently $\neg(p\Leftrightarrow q)$, Exercise 1.3.38).

\sol{1.E.14} The column equals the column of $q$: the formula $q$ (or $(p\vee\neg p)\wedge q$).

\sol{1.E.15} The column is $\checkmark$ exactly in the rows where $p$ and $q$ have the same value, regardless of $r$: $p\Leftrightarrow q$.

\sol{1.E.16} The column is $\checkmark$ exactly when $r$ is $\checkmark$ and $(p,q)\ne(\checkmark,\times)$: $r\wedge(p\Rightarrow q)$.

\sol{1.E.17} First note $p\uparrow q\equiv\neg(p\wedge q)$: if $\neg(p\wedge q)$ then by de Morgan $\neg p$ or $\neg q$, and rule (i) or (ii) gives $p\uparrow q$; conversely if $p\uparrow q$, assume $p\wedge q$, so $p$ and $q$, and rule (iii) yields $\bot$; hence $\neg(p\wedge q)$. (a) $p\uparrow p\equiv\neg(p\wedge p)\equiv\neg p$; so $(p\uparrow p)\uparrow(p\uparrow p)\equiv\neg(p\uparrow p)\equiv\neg\neg p\equiv p$. (b) $(p\uparrow p)\uparrow(q\uparrow q)\equiv\neg(\neg p\wedge\neg q)\equiv p\vee q$ (de Morgan); $(p\uparrow q)\uparrow(p\uparrow q)\equiv\neg(p\uparrow q)\equiv\neg\neg(p\wedge q)\equiv p\wedge q$. (c) $p\Rightarrow q\equiv\neg p\vee q\equiv\neg(p\wedge\neg q)\equiv p\uparrow(\neg q)\equiv p\uparrow(q\uparrow q)$.

\sol{1.E.19} $\neg\exists!x\in X,\ p(x)\equiv\neg[(\exists x,\ p(x))\wedge(\forall a,\forall b,\ (p(a)\wedge p(b))\Rightarrow a=b)]\equiv(\forall x\in X,\ \neg p(x))\vee(\exists a\in X,\ \exists b\in X,\ (p(a)\wedge p(b)\wedge a\ne b))$. \emph{Strategy:} either show that no $x$ satisfies $p$, or exhibit two distinct elements that do. \emph{Application.} Let $a\in\R$ with $a\ne-1$ and $q(x)=x^4-2ax^2+a^2-1=(x^2-a)^2-1=(x^2-a-1)(x^2-a+1)$. If $a>-1$ then $a+1>0$ and $x=\sqrt{a+1}$, $x=-\sqrt{a+1}$ are two distinct real roots (distinct because $a+1\ne0$). If $a<-1$ then $a+1<0$ and $a-1<0$, so neither $x^2=a+1$ nor $x^2=a-1$ has a real solution: no real root at all. In both cases there is not a unique real root.

\sol{1.E.20} Define $\forall!x\in X,\ p(x)$ to mean $\neg\exists!x\in X,\ \neg p(x)$. Then $\neg\forall!x,\ p(x)\equiv\neg\neg\exists!x,\ \neg p(x)\equiv\exists!x,\ \neg p(x)$, and $\neg\exists!x,\ p(x)\equiv\neg\exists!x,\ \neg\neg p(x)\equiv\forall!x,\ \neg p(x)$, which are de Morgan's laws with $\forall!,\exists!$. (It says: ``it is not the case that exactly one $x$ fails $p$''.)

\sol{1.E.21} True: Theorem 1.3.19.

\sol{1.E.22} False: Exercise 1.E.12 ($p\Rightarrow q\not\equiv q\Rightarrow p$).

\sol{1.E.23} True: by de Morgan and double negation, $p\wedge q\equiv\neg(\neg p\vee\neg q)$; replacing every conjunction this way (innermost first) gives an equivalent formula with only disjunctions and negations.

\sol{1.E.24} False: $p\wedge q$ is a counterexample. A formula built from variables with disjunctions only is true as soon as one of its variables is true; if it contains $p$ it is true when $p$ is true and $q$ false, if it contains only $q$ it is true when $q$ is true and $p$ false, and if it contains neither it does not depend on $p,q$ at all; in every case it differs from $p\wedge q$ somewhere.

\sol{1.E.25} False: for $p$ false and $r$ true, $p\wedge(q\vee r)$ is false while $(p\wedge q)\vee r$ is true.

\sol{1.E.26} False: for $p$ false, $q$ true, $r$ false, $p\vee(q\vee r)$ is true while $(p\vee q)\wedge(p\vee r)$ is $\checkmark\wedge\times=\times$.

\sol{1.E.27} False: the correct negation is $\exists x\in X,\ \exists y\in Y,\ \neg p(x,y)$. Counterexample: $X=Y=\{0\}$ with $p(0,0)$ true; then $\neg\forall x\forall y\,p$ is false but $\exists x\forall y\,p$ is true.

\sol{1.E.31} Sometimes: if $q$ is true, $p\Rightarrow q$ is true; if $q$ is false (and $p$ true), it is false.

\sol{1.E.32} Always: assume $p$; together with $\neg p$ (i.e.\ $p$ false) this is a contradiction, so $q$ follows by the principle of explosion; hence $p\Rightarrow q$ (``vacuously true'').

\sol{1.E.33} Always: assume $\forall x\forall y\,p(x,y)$; let $y\in Y$ and $x\in X$; then $p(x,y)$ by ($\forall$E) twice; so $\forall y\forall x\,p(x,y)$; and symmetrically.

\sol{1.E.34} Sometimes. Theorem 1.2.40 gives $\exists\forall\Rightarrow\forall\exists$ always. The converse can hold (e.g.\ $p(x,y)$: ``$y=y$'', both true) and can fail (Exercise 1.2.39: $p(x,y)$: ``$x+y$ even'' on $\Z$).

\sol{1.E.35} Never. By de Morgan's laws for quantifiers, $\exists x\exists y\,\neg p(x,y)\equiv\neg\forall x\forall y\,p(x,y)$, so the second formula is the negation of the first; a formula and its negation always have opposite truth values, so they are never equivalent.

\hsection{Chapter 2}
\sol{2.1.1} If $R\in R$, then by definition of $R$ we have $R\in S$ and $R\notin R$: contradiction. If $R\notin R$, then $R$ is a set, so $R\in S$, and $R\notin R$, so $R$ satisfies the defining property of $R$, giving $R\in R$: contradiction. Neither truth value is possible: a paradox.

\sol{2.1.3} ``For all $x\in X$'' must say nothing about objects outside $X$, and an implication with false hypothesis is true, so $x\notin X$ never produces a counterexample. ``There exists $x\in X$'' must produce a witness that is in $X$ \emph{and} satisfies $p$; with $\Rightarrow$ instead, any object outside $X$ would be a witness. De Morgan: $\neg\forall x,\ (x\in X\Rightarrow p(x))\equiv\exists x,\ \neg(x\in X\Rightarrow p(x))\equiv\exists x,\ (x\in X\wedge\neg p(x))$, which is $\exists x\in X,\ \neg p(x)$; and $\neg\exists x,\ (x\in X\wedge p(x))\equiv\forall x,\ \neg(x\in X\wedge p(x))\equiv\forall x,\ (x\in X\Rightarrow\neg p(x))$ (Exercise 1.3.28), which is $\forall x\in X,\ \neg p(x)$.

\sol{2.1.5} With Definition 2.1.4, $R\in R\Leftrightarrow R\in\mathcal U\wedge R\in S\wedge R\notin R$. If $R\in\mathcal U$ we would get the paradox of 2.1.1 again, so $R\notin\mathcal U$. Then $R\in R$ is simply false (the right-hand side has the false conjunct $R\in\mathcal U$), and $R\notin R$ no longer implies $R\in R$. No contradiction remains: $R$, like $\mathcal U$ itself, is a collection that is not an element of the universe.

\sol{2.1.13} Two sets should be equal when they have exactly the same elements, regardless of how they are described: this is Axiom 2.1.14.

\sol{2.1.28} $\varnothing$ has no elements; $\{\varnothing\}$ has one element, namely $\varnothing$; $\{\{\varnothing\}\}$ has one element, namely $\{\varnothing\}$, which is not $\varnothing$. (a) $\varnothing\in\varnothing$: false. (b) $\varnothing\subseteq\varnothing$: true. (c) $\varnothing\in\{\varnothing\}$: true. (d) $\varnothing\in\{\{\varnothing\}\}$: false (its only element is $\{\varnothing\}\ne\varnothing$). (e) $\{\varnothing\}\subseteq\{\varnothing\}$: true. (f) $\{\varnothing\}\subseteq\{\{\varnothing\}\}$: false ($\varnothing$ is in the left set but not in the right). More: $\{\varnothing\}\in\{\{\varnothing\}\}$ true; $\{\varnothing\}\in\{\varnothing\}$ false; $\varnothing\subseteq\{\{\varnothing\}\}$ true; $\{\{\varnothing\}\}\subseteq\{\varnothing\}$ false.

\sol{2.1.33} $\PS(\varnothing)=\{\varnothing\}$; $\PS(\PS(\varnothing))=\{\varnothing,\{\varnothing\}\}$; $\PS(\PS(\PS(\varnothing)))=\{\varnothing,\ \{\varnothing\},\ \{\{\varnothing\}\},\ \{\varnothing,\{\varnothing\}\}\}$.

\sol{2.2.7} (a) $A\cap(0,\infty)$. (b) $A\cap\Z$.

\sol{2.2.12} ($\Leftarrow$) Suppose $b\le c$, and let $x\in(a,b)\cap(c,d)$. Then $x<b\le c<x$, contradiction; so the intervals are disjoint. The case $d\le a$ is symmetric. ($\Rightarrow$) By contraposition: suppose $b>c$ and $d>a$. Then $\max(a,c)<\min(b,d)$ (each of $a,c$ is less than each of $b,d$), so $x=\frac{\max(a,c)+\min(b,d)}2$ satisfies $a<x<b$ and $c<x<d$, and the intervals are not disjoint.

\sol{2.2.15} $A\cap[0,\infty)\cap\Q$.

\sol{2.2.20} $x\in\bigcap_{n\in\Z}A_n$ iff $x\in\R$ and $x\ne n$ for every integer $n$, i.e.\ iff $x$ is a real number that is not an integer. So the intersection is the set of non-integer real numbers ($\R\setminus\Z$ in the notation of 2.2.38).

\sol{2.2.31} $x^4-5x^2+4=(x^2-1)(x^2-4)$ is positive iff both factors are positive or both are negative, i.e.\ iff $x^2>4$ (then also $x^2>1$) or $x^2<1$ (then also $x^2<4$). Hence $X=(-\infty,-2)\cup(-1,1)\cup(2,\infty)$.

\sol{2.2.32} For sets $X,A_0,\dots,A_{n-1}$: $X\cap(A_0\cup\dots\cup A_{n-1})=(X\cap A_0)\cup\dots\cup(X\cap A_{n-1})$ and $X\cup(A_0\cap\dots\cap A_{n-1})=(X\cup A_0)\cap\dots\cap(X\cup A_{n-1})$. \emph{Proof of the first:} $a\in X\cap(A_0\cup\dots\cup A_{n-1})\Leftrightarrow a\in X\wedge\exists i\in[n],\ a\in A_i\Leftrightarrow\exists i\in[n],\ (a\in X\wedge a\in A_i)\Leftrightarrow\exists i\in[n],\ a\in X\cap A_i\Leftrightarrow a\in(X\cap A_0)\cup\dots\cup(X\cap A_{n-1})$. \emph{Second:} $a\in X\cup(A_0\cap\dots\cap A_{n-1})\Leftrightarrow a\in X\vee\forall i\in[n],\ a\in A_i\Leftrightarrow\forall i\in[n],\ (a\in X\vee a\in A_i)\Leftrightarrow a\in(X\cup A_0)\cap\dots\cap(X\cup A_{n-1})$; the middle step is the law $p\vee\forall i\,q(i)\equiv\forall i\,(p\vee q(i))$ when $p$ does not involve $i$ (prove it by cases on $p$).

\sol{2.2.42} (a) $A=\R\setminus\Z$. (b) $A=\bigcup_{n\in\Z}(n,n+1)$: a non-integer real $x$ lies in $(\lfloor x\rfloor,\lfloor x\rfloor+1)$, and every element of some $(n,n+1)$ is a real number strictly between consecutive integers, hence not an integer.

\sol{2.2.43} For any $a$: $a\in X\setminus(X\setminus A)\Leftrightarrow a\in X\wedge\neg(a\in X\wedge a\notin A)\Leftrightarrow a\in X\wedge(a\notin X\vee a\in A)$ \reason{de Morgan} $\Leftrightarrow(a\in X\wedge a\notin X)\vee(a\in X\wedge a\in A)$ \reason{distributivity, Example 1.3.3} $\Leftrightarrow a\in X\wedge a\in A$ \reason{the first disjunct is a contradiction} $\Leftrightarrow a\in X\cap A$. \emph{Deduction:} by Proposition 2.2.8 (with $A$ in the role of $X$ and $X$ in the role of $Y$), $A\subseteq X$ iff $A\cap X=A$, i.e.\ iff $X\cap A=A$, i.e.\ iff $X\setminus(X\setminus A)=A$.

\sol{2.2.55} A typical element of $\R^3=\prod_{i\in[3]}\R=\R\times\R\times\R$ is a triple $(x,y,z)$ of reals, so $\R^3=\R\times\R\times\R$ (Notation 2.2.54 defines $\R^3$ as this product). A typical element of $\R^2\times\R$ is a \emph{pair} $((x,y),z)$ whose first term is a pair, and $(\R\times\R)\times\R=\R^2\times\R$ since $\R\times\R=\R^2$. A typical element of $\R\times\R^2=\R\times(\R\times\R)$ is a pair $(x,(y,z))$. So the six sets fall into three equal groups: $\R^3=\R\times\R\times\R$; $\R^2\times\R=(\R\times\R)\times\R$; $\R\times\R^2=\R\times(\R\times\R)$. Sets from different groups are not equal: a $3$-tuple is not a $2$-tuple (different lengths, Definition 2.2.46), and $((x,y),z)\ne(x,(y,z))$ since their first terms $(x,y)$ and $x$ differ. They are all equivalent in the sense that there are one-to-one correspondences between them, e.g.\ $((x,y),z)\mapsto(x,y,z)$: this is the notion of \emph{bijection} of Section 3.2 (compare Exercise 3.2.47).

\sol{2.E.1} (a) $n^2<20$ for an integer $n$ iff $-4\le n\le4$ (as $4^2=16<20<25=5^2$): $\{-4,-3,-2,-1,0,1,2,3,4\}$. (b) $\{3,7,11,15,19,\dots\}$. (c) An odd multiple of six is $6m$ with $m$ odd: $\{6(2k+1)\mid k\in\Z\}=\{12k+6\mid k\in\Z\}$. (d) $\{n^2+1\mid n\in\N\}$.

\sol{2.E.2} Take $X_n=\{k\in\N\mid k\ge n\}$. Then $X_{n+1}\subseteq X_n$ (if $k\ge n+1$ then $k\ge n$) and $n\in X_n\setminus X_{n+1}$, so $X_{n+1}\subsetneq X_n$. No $X_n$ can be empty, whatever the choice of sets: if $X_n=\varnothing$ then $X_{n+1}\subsetneq\varnothing$, but a proper subset $U\subsetneq\varnothing$ would satisfy $U\subseteq\varnothing$, hence $U=\varnothing$ (an element of $U$ would be an element of $\varnothing$), contradicting $U\ne\varnothing$.

\sol{2.E.3} The set has the two elements $a=\varnothing$ and $b=\{\varnothing,\{\varnothing\}\}$, so its power set has the four subsets $\varnothing$, $\{a\}$, $\{b\}$, $\{a,b\}$:
$\PS(\{\varnothing,\{\varnothing,\{\varnothing\}\}\})=\{\varnothing,\ \{\varnothing\},\ \{\{\varnothing,\{\varnothing\}\}\},\ \{\varnothing,\{\varnothing,\{\varnothing\}\}\}\}$.

\sol{2.E.7} Let $A\in F$; we show $A\in\PS(F)$, i.e.\ $A\subseteq F$. Let $x\in A$. By the hypothesis applied to $A\in F$ and $x\in A$, $x\in F$. So $A\subseteq F$, and $F\subseteq\PS(F)$.

\sol{2.E.10} By Exercise 2.E.9, $X\symdiff X=(X\setminus X)\cup(X\setminus X)=\varnothing\cup\varnothing=\varnothing$ (an element of $X\setminus X$ would be both in and not in $X$). And $X\symdiff\varnothing=(X\setminus\varnothing)\cup(\varnothing\setminus X)=X\cup\varnothing=X$.

\sol{2.E.11} ($\Rightarrow$) If $X=Y$ then $X\symdiff Y=X\symdiff X=\varnothing$ by 2.E.10. ($\Leftarrow$) Suppose $X\symdiff Y=\varnothing$. Let $a\in X$. If $a\notin Y$ then $a\in X\setminus Y\subseteq X\symdiff Y=\varnothing$, impossible; so $a\in Y$. Thus $X\subseteq Y$, and symmetrically $Y\subseteq X$. Hence $X=Y$.

\sol{2.E.12} By 2.E.9, $X\symdiff Y=(X\cup Y)\setminus(X\cap Y)$. ($\Rightarrow$) If $X\cap Y=\varnothing$ then $X\symdiff Y=(X\cup Y)\setminus\varnothing=X\cup Y$. ($\Leftarrow$) Suppose $X\symdiff Y=X\cup Y$ and suppose $a\in X\cap Y$. Then $a\in X\cup Y=X\symdiff Y$, so $a\in X$ or $a\in Y$ but not both; yet $a$ is in both. Contradiction, so $X\cap Y=\varnothing$.

\sol{2.E.13} Let $a$ be arbitrary; $p$: $a\in X$, $q$: $a\in Y$, $r$: $a\in Z$. Write $s\oplus t$ for ``$s$ or $t$ but not both''. Then $a\in X\symdiff(Y\symdiff Z)$ is $p\oplus(q\oplus r)$ and $a\in(X\symdiff Y)\symdiff Z$ is $(p\oplus q)\oplus r$:
\[
\begin{array}{ccc|c|c|c|c}
p&q&r&q\oplus r&p\oplus(q\oplus r)&p\oplus q&(p\oplus q)\oplus r\\\hline
\checkmark&\checkmark&\checkmark&\times&\checkmark&\times&\checkmark\\
\checkmark&\checkmark&\times&\checkmark&\times&\times&\times\\
\checkmark&\times&\checkmark&\checkmark&\times&\checkmark&\times\\
\checkmark&\times&\times&\times&\checkmark&\checkmark&\checkmark\\
\times&\checkmark&\checkmark&\times&\times&\checkmark&\times\\
\times&\checkmark&\times&\checkmark&\checkmark&\checkmark&\checkmark\\
\times&\times&\checkmark&\checkmark&\checkmark&\times&\checkmark\\
\times&\times&\times&\times&\times&\times&\times
\end{array}
\]
The two columns agree (both say: $a$ lies in an odd number of the sets $X,Y,Z$), so the sets are equal by extensionality.

\sol{2.E.14} With $p,q,r$ as above: $a\in X\cap(Y\symdiff Z)$ is $p\wedge(q\oplus r)$ and $a\in(X\cap Y)\symdiff(X\cap Z)$ is $(p\wedge q)\oplus(p\wedge r)$. If $p$ is false both are false. If $p$ is true, $p\wedge q\equiv q$ and $p\wedge r\equiv r$, so both reduce to $q\oplus r$. Hence they agree for all $a$, and the sets are equal.

\sol{2.E.17} ($\Rightarrow$) Let $a\in U$ and take $\delta>0$ with $(a-\delta,a+\delta)\subseteq U$; then $u=a-\delta$, $v=a+\delta$ satisfy $u<a<v$ and $(u,v)\subseteq U$. ($\Leftarrow$) Let $a\in U$ and take $u<a<v$ with $(u,v)\subseteq U$. Put $\delta=\min(a-u,v-a)>0$. If $a-\delta<x<a+\delta$ then $u\le a-\delta<x<a+\delta\le v$, so $x\in(u,v)\subseteq U$. Hence $(a-\delta,a+\delta)\subseteq U$.

\sol{2.E.19} ($\Rightarrow$) Suppose $U$ is open. For each $a\in U$ choose $\delta_a>0$ with $(a-\delta_a,a+\delta_a)\subseteq U$ (Strategy 3.2.35). Take $I=U$ and, for $i\in U$, $a_i=i-\delta_i$, $b_i=i+\delta_i$. Then $U=\bigcup_{i\in U}(a_i,b_i)$: each $a\in U$ lies in its own interval $(a_a,b_a)$, and each interval is contained in $U$. ($\Leftarrow$) Suppose $U=\bigcup_{i\in I}(a_i,b_i)$ and let $a\in U$. Then $a\in(a_i,b_i)$ for some $i\in I$, and $(a_i,b_i)\subseteq U$. So $u=a_i$, $v=b_i$ satisfy the condition of Exercise 2.E.17, and $U$ is open.

\sol{2.E.20} ($\subseteq$) Let $x\in\bigcap_{i\in\N}A_i$ and let $j\in\N$. By (ii) there is $i\ge j$ with $A_i\subseteq B_j$; since $x\in A_i$, $x\in B_j$. As $j$ was arbitrary, $x\in\bigcap_jB_j$. ($\supseteq$) Let $x\in\bigcap_jB_j$ and let $i\in\N$. By (i) there is $j\ge i$ with $B_j\subseteq A_i$; since $x\in B_j$, $x\in A_i$. So $x\in\bigcap_iA_i$. (The conditions $j\ge i$, $i\ge j$ are not even needed.)

\sol{2.E.21} False. $E=\{0\}$: $\neg\forall x,\ x\in E$ is true (e.g.\ $1\notin E$) but $E\ne\varnothing$.

\sol{2.E.22} True. $\forall x,\ \neg x\in E$ says $E$ has no elements, i.e.\ $E$ is empty (Definition 2.1.19), so $E=\varnothing$ by Theorem 2.1.25.

\sol{2.E.23} True: both sets have exactly the elements $1,2,3$ (Axiom 2.1.14; order and repetition are irrelevant).

\sol{2.E.24} False: $\varnothing$ has no elements, so $\varnothing\in\varnothing$ is false.

\sol{2.E.25} True: $\{\varnothing\}$ has the single element $\varnothing$.

\sol{2.E.26} False: the only element of $\{\{\varnothing\}\}$ is $\{\varnothing\}$, and $\{\varnothing\}\ne\varnothing$ (the former has an element, the latter does not).

\sol{2.E.27} Always: both sets have exactly the elements $a$ and $b$, so they are equal by Axiom 2.1.14.

\sol{2.E.29} Sometimes. $X=Y=\varnothing$: any $a$ satisfies $a\in X\Leftrightarrow a\in Y$ (both false) and $X=Y$. $X=\{0\}$, $Y=\{1\}$, $a=2$: $a\in X\Leftrightarrow a\in Y$ holds (both false) but $X\ne Y$.

\sol{2.E.30} Never. If $X=Y$ then $a\in X$ gives $a\in Y$, contradicting $a\notin Y$.

\sol{2.E.31} Sometimes. $X=\{0\}$, $Y=\{1\}$: $X\ne Y$ and every $a\in X$ (only $a=0$) has $a\notin Y$. $X=\{0\}$, $Y=\{0,1\}$: $X\ne Y$ but $0\in X$ and $0\in Y$.

\sol{2.E.34} Sometimes. $X=\varnothing$, $Y=\{\varnothing\}$: $X\in Y$ and $X\subseteq Y$. $X=\{\varnothing\}$, $Y=\{\{\varnothing\}\}$: $X\in Y$, but $\varnothing\in X$ and $\varnothing\notin Y$ (the only element of $Y$ is $\{\varnothing\}\ne\varnothing$), so $X\nsubseteq Y$.

\sol{2.E.35} Always: $\varnothing\subseteq X$ (Exercise 2.1.27), so $\varnothing\in\PS(X)$.

\sol{2.E.36} Sometimes. $\{\varnothing\}\in\PS(X)$ iff $\{\varnothing\}\subseteq X$ iff $\varnothing\in X$. True for $X=\{\varnothing\}$; false for $X=\varnothing$ (then $\PS(X)=\{\varnothing\}$ and $\{\varnothing\}\ne\varnothing$).

\sol{2.E.37} Always: $\varnothing\subseteq X$ and $X\subseteq X$, so both elements of $\{\varnothing,X\}$ are in $\PS(X)$.

\sol{2.E.39} Sometimes. $X=Y=\varnothing$: $X\setminus Y=\varnothing$ and $X=Y$. $X=\varnothing$, $Y=\{0\}$: $X\setminus Y=\varnothing$ but $X\ne Y$. (In general $X\setminus Y=\varnothing$ iff $X\subseteq Y$.)

\sol{2.E.40} Always. First, $Y=\varnothing$: if $y\in Y$ then $y\in Z=X\setminus Y$, so $y\notin Y$, a contradiction. Hence $Z=X\setminus\varnothing=X$, and $Y\cup Z=\varnothing\cup X=X$.

\sol{2.E.42} Sometimes. $X=\{\varnothing\}$, $A=\varnothing$: $A\subseteq X$ and $A\in X$. $X=\{0\}$, $A=\{0\}$: $A\subseteq X$ but $A\notin X$ (the only element of $X$ is $0\ne\{0\}$).

\hsection{Chapter 3}
\sol{3.1.3} Only (a) has the form $\forall x\in X,\ \exists!\,y\in Y,\ p(x,y)$: ``for each real number $a$, the equation $x^2+2ax+a^2=0$ has exactly one real solution $x$'' is $\forall a\in\R,\ \exists!\,x\in\R,\ x^2+2ax+a^2=0$. Since $x^2+2ax+a^2=(x+a)^2$, the unique solution is $x=-a$, so the function is $\R\to\R$, $a\mapsto-a$. Statement (b) begins with $\exists!\,a\in\R$ and (c) with $\exists!\,n\in\N$, so neither has the required form.

\sol{3.1.22} By Definition 3.1.17, $g\circ f$ requires the codomain of $f$ to equal the domain of $g$, and $Y\ne Z$. Let $i:Y\to Z$ be the inclusion, $i(y)=y$ (well-defined since $Y\subseteq Z$). Then $h=g\circ i\circ f=(g\circ i)\circ f$: for $x\in X$, $(g\circ(i\circ f))(x)=g(i(f(x)))=g(f(x))=h(x)$, and the domains and codomains agree ($X\to W$).

\sol{3.1.28} Let $x\in X$. \emph{(b)} If $x\in A$ and $x\in B$: right side $1+1-1=1$, and $x\in A\cup B$ so left side $1$. If $x\in A$, $x\notin B$: $1+0-0=1$ and $x\in A\cup B$. If $x\notin A$, $x\in B$: symmetric. If $x\notin A$, $x\notin B$: $0+0-0=0$ and $x\notin A\cup B$. In all cases $\chi_{A\cup B}(x)=\chi_A(x)+\chi_B(x)-\chi_A(x)\chi_B(x)$. \emph{(c)} If $x\in A$ then $x\notin X\setminus A$, so $\chi_{X\setminus A}(x)=0=1-1$; if $x\notin A$ then $x\in X\setminus A$ (as $x\in X$), so $\chi_{X\setminus A}(x)=1=1-0$.

\sol{3.1.30} For $x\in X$, by Theorem 3.1.27: $\chi_{X\setminus(A\cup B)}=1-\chi_{A\cup B}=1-\chi_A-\chi_B+\chi_A\chi_B=(1-\chi_A)(1-\chi_B)=\chi_{X\setminus A}\chi_{X\setminus B}=\chi_{(X\setminus A)\cap(X\setminus B)}$, so $X\setminus(A\cup B)=(X\setminus A)\cap(X\setminus B)$ by Strategy 3.1.26. Likewise $\chi_{(X\setminus A)\cup(X\setminus B)}=(1-\chi_A)+(1-\chi_B)-(1-\chi_A)(1-\chi_B)=2-\chi_A-\chi_B-(1-\chi_A-\chi_B+\chi_A\chi_B)=1-\chi_A\chi_B=1-\chi_{A\cap B}=\chi_{X\setminus(A\cap B)}$, giving $X\setminus(A\cap B)=(X\setminus A)\cup(X\setminus B)$.

\sol{3.1.34} If $y\in f[\varnothing]$ then $y=f(x)$ for some $x\in\varnothing$, which is impossible. So $f[\varnothing]$ has no elements: $f[\varnothing]=\varnothing$.

\sol{3.1.42} Let $A=f^{-1}[\{1\}]=\{x\in X\mid f(x)=1\}$. Both $f$ and $\chi_A$ have domain $X$ and codomain $\{0,1\}=[2]$. Let $x\in X$. If $f(x)=1$ then $x\in A$, so $\chi_A(x)=1=f(x)$. If $f(x)\ne1$ then $f(x)=0$ (the only other element of the codomain) and $x\notin A$, so $\chi_A(x)=0=f(x)$. Hence $f=\chi_A$ by function extensionality.

\sol{3.2.7} $f:\N\to\Z$, $f(n)=n^2$: injective. If $m,n\in\N$ and $m^2=n^2$ then $(m-n)(m+n)=0$; if $m+n=0$ then $m=n=0$; otherwise $m-n=0$. Either way $m=n$. \ $g:\Z\to\N$, $g(n)=n^2$: not injective, $g(-1)=1=g(1)$. \ $h(x,y,z)=2^x3^y5^z$: injective, by uniqueness of prime factorisation (the fundamental theorem of arithmetic, Theorem 7.2.12, which the book also uses in Example 3.2.24): if $2^x3^y5^z=2^{x'}3^{y'}5^{z'}$ then the exponents of each prime agree, so $(x,y,z)=(x',y',z')$.

\sol{3.2.13} Take $V=f[X]$ and $g:X\to f[X]$, $g(x)=f(x)$ (well-defined since $f(x)\in f[X]$). For $y\in f[X]$ there is $x\in X$ with $f(x)=y$ (Definition 3.1.31), so $g(x)=y$: $g$ is surjective.

\sol{3.2.15} Take $Z=f[X]$, $p:X\to f[X]$, $p(x)=f(x)$ (surjective by Exercise 3.2.13) and $i:f[X]\to Y$, $i(y)=y$ (injective: $i(y)=i(y')$ means $y=y'$). Then $(i\circ p)(x)=i(f(x))=f(x)$ for all $x\in X$, so $f=i\circ p$.

\sol{3.2.21} ($\Rightarrow$) Suppose $f$ injective and let $y\in f[X]$. Existence: $y=f(x)$ for some $x\in X$ by definition of $f[X]$. Uniqueness: if $y=f(x)=f(x')$ then $x=x'$ by injectivity. ($\Leftarrow$) Suppose the condition holds and let $a,b\in X$ with $f(a)=f(b)=:y$. Then $y\in f[X]$, and both $a$ and $b$ are elements $x$ with $y=f(x)$; by uniqueness $a=b$.

\sol{3.2.22} Injective: if $2^m3^n=2^{m'}3^{n'}$ then $m=m'$ and $n=n'$ by uniqueness of prime factorisation (Theorem 7.2.12). (Directly: if $m<m'$, dividing by $2^m$ gives $3^n=2^{m'-m}3^{n'}$ with left side odd and right side even, impossible; so $m=m'$, then $3^n=3^{n'}$ gives $n=n'$.) Decoding: $1=2^03^0\mapsto(0,0)$; $24=2^3\cdot3\mapsto(3,1)$; $7776=6^5=2^53^5\mapsto(5,5)$; $59049=3^{10}\mapsto(0,10)$; $396718580736=2^{10}\cdot3^{18}\mapsto(10,18)$.

\sol{3.2.28} Let $x\in X$. Then $x=g(f(x))$, so $x$ is the value of $g$ at $f(x)\in Y$. Hence $g$ is surjective.

\sol{3.2.31} Let $g:Y\to X$ with $f\circ g=\id_Y$ and let $y\in Y$. Then $x=g(y)\in X$ satisfies $f(x)=f(g(y))=y$. So $f$ is surjective.

\sol{3.2.39(c)} Every $k\in\N$ has $k+1\ge1$, and every positive integer can be written uniquely as $2^m q$ with $m\in\N$ and $q$ odd (fundamental theorem of arithmetic: $m$ is the exponent of $2$ and $q$ the product of the odd prime powers; directly: if $2^mq=2^{m'}q'$ with $q,q'$ odd and $m<m'$ then $q=2^{m'-m}q'$ is even, contradiction, so $m=m'$ and $q=q'$). Write $q=2n+1$ and define $h^{-1}:\N\to\N\times\N$ by $h^{-1}(k)=(m,n)$ where $k+1=2^m(2n+1)$. Then $h^{-1}(h(m,n))=h^{-1}(2^m(2n+1)-1)=(m,n)$ by uniqueness, and $h(h^{-1}(k))=2^m(2n+1)-1=(k+1)-1=k$. So $h^{-1}$ is a two-sided inverse.

\sol{3.E.1} (a) Yes: $x+y=1\Leftrightarrow y=1-x$; $f(x)=1-x$. (b) No: for $x=0$ both $y=1$ and $y=-1$ satisfy the equation, so no single value $f(0)$ works (and for $x=2$ no $y$ at all). (c) No: for $x=0$ the equation ``$x=0$'' is true for every $y$, but $y=f(0)$ holds for only one $y$. (d) Yes: $f(x)=0$ for all $x$. (e) Yes: $1+x^2\ne0$, so $(1+x^2)y=1\Leftrightarrow y=\frac1{1+x^2}$. (f) No: for $x=1$ every $y$ satisfies $(1-1)y=0$.

\sol{3.E.3} Let $f:X\to\varnothing$ be a function. If $X$ had an element $a$, then $f(a)\in\varnothing$, impossible; so $X=\varnothing$. Any two functions $\varnothing\to\varnothing$ have the same domain and codomain and (vacuously) the same values, so they are equal by Axiom 3.1.4. Hence there is exactly one function with empty codomain, the empty function $\varnothing\to\varnothing$; its domain is $\varnothing$.

\sol{3.E.5} $f(-x)=(-x)^n=(-1)^nx^n$. If $n$ is even, $(-1)^n=1$ so $f(-x)=f(x)$: $f$ is even; if $n$ is odd, $(-1)^n=-1$ so $f(-x)=-f(x)$: $f$ is odd. Conversely, if $f$ is even then $f(-1)=f(1)$ gives $(-1)^n=1$, so $n$ is even (for odd $n$, $(-1)^n=-1\ne1$); if $f$ is odd then $(-1)^n=-1$, so $n$ is odd.

\sol{3.E.6} If $f$ is even and odd then for all $x$, $f(x)=f(-x)=-f(x)$, so $2f(x)=0$, $f(x)=0$. Thus $f$ is the zero function, which is indeed even and odd ($0=0=-0$). So exactly one such function exists.

\sol{3.E.7} Write $-U=\{-u\mid u\in U\}$; note $x\in-U\Leftrightarrow-x\in U$ (if $x=-u$ with $u\in U$ then $-x=u\in U$; conversely $x=-(-x)$). (a) $\chi_U$ even $\Leftrightarrow\forall x,\ \chi_U(-x)=\chi_U(x)\Leftrightarrow\forall x,\ (-x\in U\Leftrightarrow x\in U)\Leftrightarrow\forall x,\ (x\in-U\Leftrightarrow x\in U)\Leftrightarrow-U=U$. (b) $\chi_U$ odd $\Leftrightarrow\forall x,\ \chi_U(-x)=-\chi_U(x)\Leftrightarrow\forall x,\ (-x\in U\Leftrightarrow x\notin U)\Leftrightarrow\forall x,\ (x\in-U\Leftrightarrow x\in\R\setminus U)\Leftrightarrow-U=\R\setminus U$.

\sol{3.E.8} Put $g(x)=\frac{f(x)+f(-x)}2$ and $h(x)=\frac{f(x)-f(-x)}2$. Then $g(-x)=g(x)$, $h(-x)=-h(x)$ and $g(x)+h(x)=f(x)$. Uniqueness: if also $f=g'+h'$ with $g'$ even, $h'$ odd, then $g-g'=h'-h$; the left side is even and the right side is odd, so this function is both, hence zero by Exercise 3.E.6. So $g=g'$ and $h=h'$.

\sol{3.E.9} For $n=0$: $[0]=\varnothing$, and $\theta_0:\varnothing\to\varnothing$ is the empty function, which is $\id_\varnothing$. For $n\ge1$: first, $\theta_1=\id_{[1]}$, since $[1]=\{0\}$ has only one function to itself (3.E.39). Fix $k\in[n]$ and let $f:[1]\to[n]$, $f(0)=k$. Then $f\circ\theta_1=\theta_n\circ f$ gives $k=f(0)=f(\theta_1(0))=\theta_n(f(0))=\theta_n(k)$. As $k$ was arbitrary, $\theta_n=\id_{[n]}$.

\sol{3.E.10} $\chi_{U\symdiff V}=\chi_U+\chi_V-2\chi_U\chi_V$. Check by cases: $x$ in both: $1+1-2=0$ and $x\notin U\symdiff V$; in exactly one: $1+0-0=1$ and $x\in U\symdiff V$; in neither: $0$. (Equivalently $(\chi_U-\chi_V)^2$.)

\sol{3.E.11} On $\R\setminus\{0\}=(-\infty,0)\cup(0,\infty)$, a function $F$ with $F'(x)=\frac1x$ satisfies $(F(x)-\log|x|)'=0$ separately on each of the two intervals, so $F-\log|x|$ is constant on $(-\infty,0)$ (say $a$) and constant on $(0,\infty)$ (say $b$), and these constants are independent. Hence $F(x)=\log|x|+a\chi_U(x)+b\chi_V(x)$ with $U=(-\infty,0)$, $V=(0,\infty)$; conversely every such $F$ has derivative $\frac1x$.

\sol{3.E.12} $f[\R]=[1,\infty)$. ($\subseteq$) $1+x^2\ge1$ so $\sqrt{1+x^2}\ge1$. ($\supseteq$) For $y\ge1$ put $x=\sqrt{y^2-1}\in\R$; then $f(x)=\sqrt{1+y^2-1}=y$.

\sol{3.E.14} $f(1)=\{0\}$, $f(2)=\{0,1\}$, and in general $f(n)=[n]=\{0,1,\dots,n-1\}$ (induction: $f(n+1)=[n]\cup\{n\}=[n+1]$). So $f[\N]=\{[n]\mid n\in\N\}=\{\varnothing,\{0\},\{0,1\},\{0,1,2\},\dots\}$.

\sol{3.E.15} $f[\R^\R]$ is the set of even functions $\R\to\R$. ($\subseteq$) $f(h)(-x)=h(|-x|)=h(|x|)=f(h)(x)$, so every $f(h)$ is even. ($\supseteq$) Let $g$ be even. Then $f(g)(x)=g(|x|)=g(x)$ for all $x$ (if $x\ge0$ then $|x|=x$; if $x<0$ then $g(|x|)=g(-x)=g(x)$), so $f(g)=g$ and $g\in f[\R^\R]$.

\sol{3.E.16} $\sqrt{1+x^2}\in(-5,5]$ iff $\sqrt{1+x^2}\le5$ (it is always $>0$) iff $1+x^2\le25$ iff $x^2\le24$ iff $-2\sqrt6\le x\le2\sqrt6$. So $f^{-1}[(-5,5]]=[-2\sqrt6,2\sqrt6]$.

\sol{3.E.18} $f^{-1}[\{\varnothing\}]=\{(A,B)\in\PS(\N)\times\PS(\N)\mid A\cap B=\varnothing\}$, the set of pairs of disjoint subsets of $\N$.

\sol{3.E.20} \emph{Reading:} $v:A\to f[X]$ should satisfy $v(a)=i(a)$ for all $a\in A$. Let $q:X\to f[X]$, $q(x)=f(x)$, and $j:f[X]\to Y$, $j(y)=y$; $j$ is injective and $j\circ q=f$. By (iii) there is a unique $u:A\to f[X]$ with $j\circ u=i$, i.e.\ $u(a)=i(a)$ for all $a\in A$. So $v=u$ is the only candidate (uniqueness). $v$ is injective: $v(a)=v(b)\Rightarrow i(a)=i(b)\Rightarrow a=b$ by (i). $v$ is surjective: let $y\in f[X]$, say $y=f(x)$; by (ii) $y=i(p(x))=v(p(x))$. So $v$ is the unique bijection with $v(a)=i(a)$.

\sol{3.E.21} This is Exercise 3.1.40 (worked in Section A).

\sol{3.E.22} (a) $z\in(g\circ f)[U]\Leftrightarrow\exists x\in U,\ z=g(f(x))$; and $z\in g[f[U]]\Leftrightarrow\exists y\in f[U],\ z=g(y)\Leftrightarrow\exists x\in U,\ z=g(f(x))$ (unfold $y\in f[U]$). (b) $x\in(g\circ f)^{-1}[W]\Leftrightarrow g(f(x))\in W\Leftrightarrow f(x)\in g^{-1}[W]\Leftrightarrow x\in f^{-1}[g^{-1}[W]]$.

\sol{3.E.23} ($\subseteq$) Let $(x,z)\in p[S]$; then $(x,z)=p(x,y,z)$ for some $(x,y,z)\in S$, so $y=f(x)$ and $z=g(y)=g(f(x))=(g\circ f)(x)$; hence $(x,z)\in\mathrm{Gr}(g\circ f)$. ($\supseteq$) Let $(x,(g\circ f)(x))\in\mathrm{Gr}(g\circ f)$. Put $y=f(x)$, $z=g(y)$; then $(x,y,z)\in S$ and $p(x,y,z)=(x,z)=(x,(g\circ f)(x))$.

\sol{3.E.26(a)} Define $f:\Z\to\N$ by $f(n)=2n$ if $n\ge0$ and $f(n)=-2n-1$ if $n<0$ (values in $\N$: $2n\ge0$; for $n\le-1$, $-2n-1\ge1$). Inverse $g:\N\to\Z$: $g(m)=\frac m2$ if $m$ even, $g(m)=-\frac{m+1}2$ if $m$ odd. Check: $g(f(n))=g(2n)=n$ for $n\ge0$; for $n<0$, $f(n)=-2n-1$ is odd and $g(f(n))=-\frac{-2n-1+1}2=n$. And $f(g(m))=2\cdot\frac m2=m$ for even $m$; for odd $m$, $g(m)=-\frac{m+1}2<0$ and $f(g(m))=-2(-\frac{m+1}2)-1=m$. So $f$ is a bijection (Strategy 3.2.41).

\sol{3.E.26(c)} Let $A=\{0,1\}\cup\{\frac1n\mid n\ge2\}\subseteq[0,1]$ and $A'=\{\frac1n\mid n\ge2\}\subseteq(0,1)$. Define $f:[0,1]\to(0,1)$ by $f(0)=\frac12$, $f(1)=\frac13$, $f(\frac1n)=\frac1{n+2}$ for $n\ge2$, and $f(x)=x$ for $x\notin A$ (then $0<x<1$). Define $g:(0,1)\to[0,1]$ by $g(\frac12)=0$, $g(\frac13)=1$, $g(\frac1n)=\frac1{n-2}$ for $n\ge4$, and $g(y)=y$ for $y\notin A'$. Then $g(f(x))=x$ and $f(g(y))=y$ in every case ($f$ maps $A$ onto $A'$ by the shift $0,1,\frac12,\frac13,\dots\mapsto\frac12,\frac13,\frac14,\frac15,\dots$, and is the identity off $A$, where $[0,1]\setminus A=(0,1)\setminus A'$). So $f$ is a bijection.

\sol{3.E.26(d)} Compose bijections (Exercise 3.2.20): $[a,b]\to[0,1]$, $x\mapsto\frac{x-a}{b-a}$ (inverse $t\mapsto a+(b-a)t$); then $[0,1]\to(0,1)$ from (c); then $(0,1)\to(c,d)$, $t\mapsto c+(d-c)t$ (inverse $y\mapsto\frac{y-c}{d-c}$).

\sol{3.E.27} Write $T(k)=\binom{k+1}2=\frac{k(k+1)}2$, so $f(a,b)=T(a+b)+b$ and $T(k+1)=T(k)+(k+1)$. \emph{Claim:} for every $n\in\N$ there is a unique $k\in\N$ with $T(k)\le n<T(k+1)$. Existence: the set $\{k\in\N\mid T(k+1)>n\}$ is non-empty ($T(n+1)\ge n+1>n$), so it has a least element $k$ (well-ordering); then $T(k)\le n$ (if $k=0$, $T(0)=0\le n$; if $k>0$, minimality gives $T(k)\le n$). Uniqueness: $T$ is strictly increasing, so if $k<k'$ then $T(k+1)\le T(k')$, and $n<T(k+1)\le T(k')\le n$ is impossible. \emph{Surjective:} given $n$, take this $k$, put $b=n-T(k)$ and $a=k-b$. Then $0\le b<T(k+1)-T(k)=k+1$, so $b\le k$ and $a\in\N$; and $f(a,b)=T(k)+b=n$. \emph{Injective:} if $f(a,b)=f(a',b')=n$, then with $k=a+b$ we have $T(k)\le n=T(k)+b\le T(k)+k<T(k+1)$, so $k$ is the unique number of the claim; likewise $a'+b'=k$. Then $b=n-T(k)=b'$ and $a=k-b=a'$.

\sol{3.E.28} Let $Y=\{x\in X\mid e(x)=x\}$. Define $f:X\to Y$ by $f(x)=e(x)$: well-defined since $e(e(x))=e(x)$ shows $e(x)\in Y$. Define $g:Y\to X$ by $g(y)=y$. Then $g(f(x))=e(x)$, so $g\circ f=e$; and for $y\in Y$, $f(g(y))=e(y)=y$, so $f\circ g=\id_Y$.

\sol{3.E.29} False: $(1,1)$ and $(1,-1)$ are both in $G$, so uniqueness fails at $x=1$ (Theorem 3.1.9).

\sol{3.E.30} True: for each $x\in\Z$ the unique $y\in\N$ with $y^2=x^2$ is $y=|x|$ (if $y^2=x^2$ then $y=\pm x$, and only one of $x,-x$ is in $\N$ unless $x=0$). So $G=\mathrm{Gr}(x\mapsto|x|)$.

\sol{3.E.31} False: the empty function $\varnothing\to\{0\}$ (Definition 3.1.13) has empty domain and non-empty codomain.

\sol{3.E.32} True: this is Exercise 3.E.3.

\sol{3.E.33} False: $f:\{0\}\to\{1\}$, $f(0)=1$, has image $\{1\}\nsubseteq\{0\}$. (The image is a subset of the \emph{codomain}.)

\sol{3.E.36} True. A left inverse $g$ of $f$ is surjective by Exercise 3.2.28. If $g$ is a right inverse of $f$ ($f\circ g=\id_Y$) then $f$ is a left inverse of $g$, so $g$ is injective by Exercise 3.2.25.

\sol{3.E.39} Always: $f(x)=0$ for all $x\in X$ is well-defined, and any function $X\to\{0\}$ must take the value $0$ everywhere, so by function extensionality it equals $f$.

\sol{3.E.40} Sometimes. If $X=\varnothing$, the empty function $\varnothing\to\varnothing$ is the unique function (3.E.3). If $X\ne\varnothing$ there is no function $X\to\varnothing$ at all (3.E.3), so ``a unique function'' is false.

\sol{3.E.42} Never. The three values $f(1),f(2),f(3)$ lie in $\{1,2\}$, so two of them coincide: if $f(1)=f(2)$ then $(f(1),1),(f(1),2)\in G'$; otherwise $\{f(1),f(2)\}=\{1,2\}$ and $f(3)=f(k)$ for some $k\in\{1,2\}$, giving $(f(3),3),(f(3),k)\in G'$ with $k\ne3$. In all cases some first coordinate occurs twice in $G'$, so $G'$ is not a graph (Theorem 3.1.9).

\sol{3.E.43} Always: pick $x\in U$; then $f(x)\in f[U]$.

\sol{3.E.44} Sometimes. $f:\R\to\R$, $f(x)=x^2$: $V=\{-1\}$ gives $f^{-1}[V]=\varnothing$; $V=\{1\}$ gives $f^{-1}[V]=\{-1,1\}\ne\varnothing$.

\sol{3.E.45} Sometimes: $f(1)=1$, $f(2)=2$ is injective; $f(1)=f(2)=1$ is not.

\sol{3.E.46} Never: $f[\{1,2\}]=\{f(1),f(2)\}$ has at most two elements, so at least one of $1,2,3$ is not a value of $f$. (Precisely: if $f(1)=f(2)$ two of the three are missed; otherwise exactly the third element is missed.)

\hsection{Chapter 4}
\sol{4.1.9} Define $m^0=1$ and $m^{n+1}=m^n\cdot m$ for all $n\in\N$. Then $2^2=2^{1+1}=2^1\cdot2=(2^0\cdot2)\cdot2=(1\cdot2)\cdot2$. Now $1\cdot2=1\cdot(1+1)=(1\cdot1)+1$ and $1\cdot1=1\cdot(0+1)=(1\cdot0)+1=0+1=1$ (using $0+1=0+s(0)=s(0+0)=s(0)=1$), so $1\cdot2=1+1=2$. Hence $2^2=2\cdot2=4$ by Proposition 4.1.8.

\sol{4.1.16} Each entry is the sum of the two above it: $1\ 6\ 15\ 20\ 15\ 6\ 1$ and $1\ 7\ 21\ 35\ 35\ 21\ 7\ 1$.

\sol{4.2.12} Elements of $\prod_{k\in[n+1]}X_k$ are $(n+1)$-tuples $(a_0,\dots,a_n)$ with $a_k\in X_k$; elements of $\big(\prod_{k\in[n]}X_k\big)\times X_n$ are pairs $((a_0,\dots,a_{n-1}),a_n)$ with the same conditions (Definition 2.2.49). Define $\varphi(a_0,\dots,a_n)=((a_0,\dots,a_{n-1}),a_n)$ and $\psi((a_0,\dots,a_{n-1}),a_n)=(a_0,\dots,a_n)$; both are well-defined, and $\psi\circ\varphi$ and $\varphi\circ\psi$ are identities by the equality rule for tuples (Definition 2.2.46), so $\varphi$ is a bijection (Strategy 3.2.41). The book asks for induction on $n$: the base case $n=0$ is $X_0\to\{()\}\times X_0$, $x\mapsto((),x)$, and the induction step is the same re-bracketing one coordinate further; the induction is what makes the ``$\dots$'' in the tuples precise.

\sol{4.2.14} For $n\ge1$; $p(n)$: $\sum_{i=0}^n(-1)^i\binom ni=0$. \emph{Base case:} $\binom10-\binom11=1-1=0$. \emph{Step:} let $n\ge1$ and suppose $p(n)$. Using $\binom{n+1}0=1$ and $\binom{n+1}i=\binom n{i-1}+\binom ni$ for $i\ge1$ (Definition 4.1.15):
\[
\sum_{i=0}^{n+1}(-1)^i\binom{n+1}i=1+\sum_{i=1}^{n+1}(-1)^i\binom n{i-1}+\sum_{i=1}^{n+1}(-1)^i\binom ni
=1-\sum_{j=0}^n(-1)^j\binom nj+\Big(\sum_{i=0}^{n+1}(-1)^i\binom ni-1\Big).
\]
The first sum is $0$ by $p(n)$; the last sum is $\sum_{i=0}^n(-1)^i\binom ni+(-1)^{n+1}\binom n{n+1}=0+0$. Total: $1-0+(0-1)=0$. The result follows by induction.

\sol{4.3.22} Suppose $\sqrt n$ is not an integer but is rational, say $\sqrt n=\frac ab$ with $a,b\in\N$, $b\ge1$. Let $m=\lfloor\sqrt n\rfloor$, so $0<\sqrt n-m<1$. Let $X=\{b\in\N\mid b\ge1,\ b\sqrt n\in\Z\}$; $X$ is non-empty since $b\sqrt n=a\in\Z$. Let $b\in X$ and put $b'=b(\sqrt n-m)=b\sqrt n-bm\in\Z$. Then $0<b'<b$, so $b'\in\N$ and $b'\ge1$; and $b'\sqrt n=bn-m\,b\sqrt n\in\Z$. So $b'\in X$ with $b'<b$: $X$ has infinite descent, contradicting the well-ordering principle (Strategy 4.3.19). Hence $\sqrt n$ is irrational. (The hint's expression $\sqrt n=\frac{a(\sqrt n-m)}{b(\sqrt n-m)}=\frac{bn-ma}{b'}$ exhibits the smaller denominator.)

\sol{4.E.1} $1+3=1+s(2)=s(1+2)=s(1+s(1))=s(s(1+1))=s(s(1+s(0)))=s(s(s(1+0)))=s(s(s(1)))=s(s(2))=s(3)=4$.

\sol{4.E.2} $0+5=s(0+4)=s(s(0+3))=s(s(s(0+2)))=s(s(s(s(0+1))))=s(s(s(s(s(0+0)))))=s(s(s(s(s(0)))))=5$, using $0+s(n)=s(0+n)$ five times and $0+0=0$.

\sol{4.E.3} $2\cdot3=2\cdot(2+1)=(2\cdot2)+2=4+2$ (Proposition 4.1.8) $=4+s(1)=s(4+1)=s(4+s(0))=s(s(4+0))=s(s(4))=s(5)=6$.

\sol{4.E.4} $0\cdot5=0\cdot(4+1)=(0\cdot4)+0=0\cdot4$, since $m+0=m$. Repeating, $0\cdot4=0\cdot3=0\cdot2=0\cdot1=(0\cdot0)+0=0\cdot0=0$.

\sol{4.E.5} With $m^{n+1}=m^n\cdot m$: $2^3=2^2\cdot2=4\cdot2$ (Exercise 4.1.9) $=4\cdot(1+1)=(4\cdot1)+4=((4\cdot0)+4)+4=(0+4)+4=4+4$ (as $0+4=4$, computed like 4.E.2) $=4+s(3)=s(4+3)=s(s(4+2))=s(s(s(4+1)))=s(s(s(s(4+0))))=s(s(s(s(4))))=8$.

\sol{4.E.6} $\exists^{=0}x\in X,\ p(x)$ means $\neg\exists x\in X,\ p(x)$. Given $\exists^{=n}$, define $\exists^{=n+1}x\in X,\ p(x)$ to mean $\exists a\in X,\ \big(p(a)\wedge\exists^{=n}x\in X,\ (p(x)\wedge x\ne a)\big)$: some $a$ satisfies $p$, and exactly $n$ elements other than $a$ do.

\sol{4.E.7} $\binom42=\binom31+\binom32$. $\binom31=\binom20+\binom21=1+\big(\binom10+\binom11\big)=1+\big(1+\big(\binom00+\binom01\big)\big)=1+(1+(1+0))=3$. $\binom32=\binom21+\binom22=2+\big(\binom11+\binom12\big)=2+\big(1+\big(\binom01+\binom02\big)\big)=2+(1+0)=3$. So $\binom42=6$.

\sol{4.E.8} The number of trailing zeros in base $b$ is the largest $r$ with $b^r\mid n$. (a) $10^r\mid41!$ iff $2^r\mid41!$ and $5^r\mid41!$. The factor $5$ occurs in $5,10,\dots,40$ (eight multiples) and once more in $25$: exponent $9$. The exponent of $2$ is larger (below), so $41!$ has $9$ trailing decimal zeros. (b) The exponent of $2$ in $41!$: $20$ multiples of $2$, $10$ of $4$, $5$ of $8$, $2$ of $16$, $1$ of $32$: $20+10+5+2+1=38$. So $41!$ has $38$ trailing zeros in binary.

\sol{4.E.9} ($\Rightarrow$) If $(N,z,s)$ satisfies Peano's axioms, apply the recursion theorem (4.1.2) with $h(n,x)=f(x)$: there is a unique $h:N\to X$ with $h(z)=a$ and $h(s(n))=f(h(n))$, i.e.\ $f\circ h=h\circ s$. ($\Leftarrow$) Assume the universal property. \emph{(iii)} Let $A\subseteq N$ with $z\in A$ and $s[A]\subseteq A$. Apply the property to $X=A$, $a=z$, $f=s|_A:A\to A$: get $h:N\to A$ with $h(z)=z$ and $s(h(n))=h(s(n))$. Composing with the inclusion $A\to N$ gives $h':N\to N$ with $h'(z)=z$, $s\circ h'=h'\circ s$; but $\id_N$ has the same properties, so $h'=\id_N$ by uniqueness (property with $X=N$, $a=z$, $f=s$). Hence $n=h(n)\in A$ for all $n$: $N\subseteq A$. \emph{(i)} Apply the property to $X=\{0,1\}$, $a=0$, $f$ the constant function $1$: $h(z)=0$ and $h(s(n))=1$ for all $n$. If $z=s(n)$ then $0=h(z)=h(s(n))=1$, contradiction. \emph{(ii)} Apply the property to $X=N\times N$, $a=(z,z)$, $f(x,y)=(y,s(y))$: get $h$ with $h(z)=(z,z)$, $h(s(n))=f(h(n))$. Writing $h(n)=(g(n),t(n))$: $t(z)=z$ and $t(s(n))=s(t(n))$, so $t=\id_N$ by uniqueness; then $h(s(n))=f(g(n),n)=(n,s(n))$ gives $g(s(n))=n$. So if $s(m)=s(n)$ then $m=g(s(m))=g(s(n))=n$.

\sol{4.E.22} Induction on $n$. \emph{Base:} $(\cos\theta+i\sin\theta)^0=1=\cos0+i\sin0$. \emph{Step:} if the formula holds for $n$ then
$(\cos\theta+i\sin\theta)^{n+1}=(\cos n\theta+i\sin n\theta)(\cos\theta+i\sin\theta)=(\cos n\theta\cos\theta-\sin n\theta\sin\theta)+i(\sin n\theta\cos\theta+\cos n\theta\sin\theta)=\cos((n+1)\theta)+i\sin((n+1)\theta)$ by the angle addition formulae.

\sol{4.E.23} Write $I_n=\int_0^{\pi/2}\sin^nx\,dx$. (a) Integrate by parts with $u=\sin^{n+1}x$, $dv=\sin x\,dx$: $I_{n+2}=\big[-\sin^{n+1}x\cos x\big]_0^{\pi/2}+(n+1)\int_0^{\pi/2}\sin^nx\cos^2x\,dx=0+(n+1)(I_n-I_{n+2})$, so $(n+2)I_{n+2}=(n+1)I_n$. (b) Induction: $I_0=\frac\pi2=\frac{\pi}{2^1}\binom00$. Step: $I_{2n+2}=\frac{2n+1}{2n+2}I_{2n}=\frac{2n+1}{2(n+1)}\cdot\frac{\pi}{2^{2n+1}}\binom{2n}n$, and $\binom{2n+2}{n+1}=\frac{(2n+2)(2n+1)}{(n+1)^2}\binom{2n}n=\frac{2(2n+1)}{n+1}\binom{2n}n$ (Theorem 4.2.15), so $\frac{\pi}{2^{2n+3}}\binom{2n+2}{n+1}=\frac{\pi(2n+1)}{2^{2n+2}(n+1)}\binom{2n}n=I_{2n+2}$. (c) $I_1=1$ and $I_{2n+3}=\frac{2n+2}{2n+3}I_{2n+1}$; by induction $I_{2n+1}=\dfrac{2^{2n}(n!)^2}{(2n+1)!}$ (base $n=0$: $1$; step: $\frac{2n+2}{2n+3}\cdot\frac{2^{2n}(n!)^2}{(2n+1)!}=\frac{2^{2n+1}(n+1)(n!)^2}{(2n+3)(2n+1)!}=\frac{2^{2n+2}((n+1)!)^2}{(2n+3)!}$, using $(2n+3)!=(2n+3)(2n+2)(2n+1)!$).

\hsection{Chapter 5}
\sol{5.1.3} For example: ``$x<y$'' on $\R$ (homogeneous); ``$x$ is a parent of $y$'' on the set of all people (homogeneous); ``the line $L$ passes through the point $P$'' from the set of lines in the plane to the set of points (not homogeneous); ``$a$ leaves remainder $r$ on division by $3$'' from $\Z$ to $[3]$ (not homogeneous).

\sol{5.1.35} ($\Rightarrow$) Let $(a,b)\in\mathrm{Gr}(R)$, so $a\mathrel Rb$. By symmetry $b\mathrel Ra$, and then by antisymmetry $a=b$; so $(a,b)=(a,a)\in\Delta_X$. ($\Leftarrow$) Suppose $\mathrm{Gr}(R)\subseteq\Delta_X$. Symmetric: if $a\mathrel Rb$ then $(a,b)\in\Delta_X$, so $a=b$, and $b\mathrel Ra$ is the true statement $a\mathrel Ra$. Antisymmetric: if $a\mathrel Rb$ (and $b\mathrel Ra$) then $(a,b)\in\Delta_X$ gives $a=b$. \emph{Deduction:} if $R$ is also reflexive then $(a,a)\in\mathrm{Gr}(R)$ for all $a$, i.e.\ $\Delta_X\subseteq\mathrm{Gr}(R)$; so $\mathrm{Gr}(R)=\Delta_X=\mathrm{Gr}(=)$ (Exercise 5.1.18), and $R$ is the equality relation by Theorem 5.1.14.

\sol{5.2.5} Reflexive: $\id_U:U\to U$ is a bijection (3.2.19), so $U\cong U$. Symmetric: if $f:U\to V$ is a bijection then $f^{-1}:V\to U$ is a bijection (3.2.45), so $V\cong U$. Transitive: if $f:U\to V$ and $g:V\to W$ are bijections then $g\circ f:U\to W$ is a bijection (3.2.20).

\sol{5.2.10} $a\equiv b\pmod0$ iff $0\mid b-a$ iff $b-a=q\cdot0=0$ for some $q$ iff $a=b$: congruence mod $0$ is equality. $a\equiv b\pmod1$ always, since $b-a=(b-a)\cdot1$ is divisible by $1$: congruence mod $1$ is the discrete relation.

\sol{5.2.14} For $r\in\{0,\dots,9\}$, $[r]_{\approx}=\{n\in\N\mid n\equiv r\pmod{10}\}=\{r,r+10,r+20,\dots\}$, the natural numbers whose last decimal digit is $r$ (Example 5.2.8). Every $n\in\N$ lies in the class of its last digit, so $\N/{\approx}=\{[0],[1],[2],[3],[4],[5],[6],[7],[8],[9]\}$.

\sol{5.2.16} By Example 5.2.15, $[a]_{\sim_f}=f^{-1}[\{f(a)\}]$, which always contains $a$. ($\Rightarrow$) If $f$ is injective and $x\in[a]$ then $f(x)=f(a)$, so $x=a$: $[a]=\{a\}$ has a unique element. ($\Leftarrow$) Suppose every class has a unique element, and let $f(a)=f(b)$. Then $b\in[a]$ and $a\in[a]$, so $a=b$ by uniqueness. Hence $f$ is injective.

\sol{5.2.20} (For $n\ne0$; for $n=0$ the statement fails, since $[0]=\varnothing$ while $\Z/0\Z=\{\{a\}\mid a\in\Z\}$ by 5.2.10.) \emph{Injective:} let $r,s\in[|n|]$ with $[r]_n=[s]_n$. Then $r\equiv s\pmod n$ (Theorem 5.2.17), so $r$ and $s$ have the same remainder on division by $n$ (5.2.9). But $0\le r,s<|n|$, so each is its own remainder ($r=0\cdot n+r$); hence $r=s$. \emph{Surjective:} let $[a]_n\in\Z/n\Z$ and let $r$ be the remainder of $a$, $a=qn+r$ with $r\in[|n|]$. Then $a-r=qn$, so $r\equiv a\pmod n$, and $[r]_n=[a]_n$ by Theorem 5.2.17; hence $[a]_n=i(r)$.

\sol{5.E.1} On $X=\{0,1,2\}$, give the graph. None: $\{(0,1),(1,2)\}$ (no $(0,0)$; no $(1,0)$; no $(0,2)$). Reflexive only: $\Delta_X\cup\{(0,1),(1,2)\}$. Symmetric only: $\{(0,1),(1,0)\}$ (no $(0,0)$; transitivity would need $(0,0)$). Transitive only: $\{(0,1)\}$ (vacuously transitive). Reflexive and symmetric only: $\Delta_X\cup\{(0,1),(1,0),(1,2),(2,1)\}$ (no $(0,2)$). Reflexive and transitive only: $\le$ on $\R$. Symmetric and transitive only: the empty relation on $X$ (vacuous; not reflexive). All three: $=$.

\sol{5.E.2} This is the ($\Rightarrow$) direction of Exercise 5.1.35 together with 5.1.18: $\mathrm{Gr}(R)\subseteq\Delta_X=\mathrm{Gr}(=)$.

\sol{5.E.3} Let $x\in X$. Since $R$ is left-total there is $y\in X$ with $x\mathrel Ry$. By symmetry $y\mathrel Rx$. By transitivity ($x\mathrel Ry$ and $y\mathrel Rx$) $x\mathrel Rx$. So $R$ is reflexive.

\sol{5.E.5} Let $S$ be the transport of $\sim$ along $f$. \emph{Symmetric, always:} if $c\mathrel Sd$ via $a,b$ ($f(a)=c$, $f(b)=d$, $a\sim b$) then $b\sim a$, so $d\mathrel Sc$ via $b,a$. \emph{Reflexive iff $f$ is surjective:} $c\mathrel Sc$ needs some $a$ with $f(a)=c$ (then $a\sim a$ works), so reflexivity on all of $Y$ holds exactly when every $c\in Y$ is a value of $f$. \emph{Transitive: not in general.} Take $X=\{0,1,2,3\}$ with $\sim$ having classes $\{0,1\}$, $\{2,3\}$, and $f(0)=c$, $f(1)=f(2)=d$, $f(3)=e$. Then $c\mathrel Sd$ (via $0\sim1$) and $d\mathrel Se$ (via $2\sim3$), but $c\mathrel Se$ would need $0\sim3$, which is false. Transitivity does hold under the extra condition ``$f(b)=f(b')\Rightarrow b\sim b'$'': if $c\mathrel Sd$ via $a\sim b$ and $d\mathrel Se$ via $b'\sim a'$ with $f(b)=f(b')=d$, then $a\sim b\sim b'\sim a'$ gives $c\mathrel Se$. \emph{Conclusion:} $S$ is an equivalence relation on $Y$ when $f$ is surjective and $f(b)=f(b')\Rightarrow b\sim b'$ (for instance when $f$ is a bijection, or when $f$ is the quotient function $q_\sim$); in general only symmetry is guaranteed.

\sol{5.E.7} Here $a\mathrel Rb$ or $b\mathrel Ra$ means $b=a+n$ or $a=b+n$, i.e.\ $b-a=\pm n$. If $x\sim_Ry$ via a chain $x=a_0,\dots,a_k=y$, then $y-x=\sum_{i\in[k]}(a_{i+1}-a_i)$ is a sum of $k$ terms each equal to $\pm n$, hence a multiple of $n$: $x\equiv y\pmod n$. Conversely if $y-x=mn$, take the chain $x,x+n,x+2n,\dots,x+mn=y$ if $m\ge0$ (each step $a_i\mathrel Ra_{i+1}$) or $x,x-n,\dots,y$ if $m<0$ (each step $a_{i+1}\mathrel Ra_i$); for $m=0$ take $k=0$. So $\sim_R$ and $\equiv\pmod n$ relate the same pairs and are equal.

\sol{5.E.8} Let $U,V\subseteq X$. Then $U\subseteq X$ and $V\subseteq X$, so the chain $U,X,V$ ($k=2$) has $U\mathrel RX$ and $V\mathrel RX$ (the second step in the reverse direction). Hence $U\sim_RV$.

\sol{5.E.9} ($\Leftarrow$) If $x=y$ take the chain of length $0$. If $\{x,y\}=\{a,b\}$, the chain $a,b$ or $b,a$ works since $a\mathrel Rb$. ($\Rightarrow$) Let $x=a_0,\dots,a_k=y$ be a chain. Each step has $a_i\mathrel Ra_{i+1}$ or $a_{i+1}\mathrel Ra_i$, and the only related pair is $(a,b)$, so every step is between $a$ and $b$. If $k=0$ then $x=y$. If $k\ge1$ then $x,y\in\{a,b\}$, so $x=y$ or $\{x,y\}=\{a,b\}$.

\sol{5.E.10} Reflexive: the chain $x$ alone ($k=0$) gives $x\sim_Rx$. Symmetric: reverse the chain ($a_k,\dots,a_0$); each step still satisfies ``$R$ in one direction or the other''. Transitive: concatenate a chain from $x$ to $y$ with a chain from $y$ to $z$.

\sol{5.E.11} If $x\mathrel Ry$ then $x,y$ is a chain with $k=1$, so $x\sim_Ry$.

\sol{5.E.12} Let $\approx$ be an equivalence relation with $x\mathrel Ry\Rightarrow x\approx y$, and let $x\sim_Ry$ via $x=a_0,\dots,a_k=y$. For each $i$, $a_i\mathrel Ra_{i+1}$ gives $a_i\approx a_{i+1}$, and $a_{i+1}\mathrel Ra_i$ gives $a_{i+1}\approx a_i$, hence $a_i\approx a_{i+1}$ by symmetry. By transitivity along the chain (Proposition 5.1.42) $a_0\approx a_k$, i.e.\ $x\approx y$; for $k=0$ use reflexivity.

\sol{5.E.13} By 5.E.11, $x\mathrel Ry\Rightarrow x\sim_Ry$. By 5.E.12 applied with $\approx=R$ (an equivalence relation containing $R$ trivially), $x\sim_Ry\Rightarrow x\mathrel Ry$. Same domain, same pairs: $\sim_R=R$ (Axiom 5.1.4).

\sol{5.E.14} (a) $\sim$ is $\sim_D$ for the function $D:C^1(\R)\to$ (functions $\R\to\R$), $D(F)=F'$, so it is an equivalence relation by Exercise 5.2.4. (b) $G\in[F]_\sim$ iff $G'=F'$ iff $(G-F)'=0$ iff $G-F$ is a constant function $c$ (a function on $\R$ with zero derivative is constant), i.e.\ $G(x)=F(x)+c$ for all $x$. (c) $\int f(x)\,dx$ ``$=F(x)+c$'' is a description of the whole equivalence class $[F]_\sim$ of one antiderivative $F$: the indefinite integral is the $\sim$-class of antiderivatives, and ``$+c$'' is exactly the content of (b).

\sol{5.E.15} False: $\le$ on $\R$ is reflexive but $0\le1$ and $1\not\le0$.

\sol{5.E.16} True: by Exercise 5.1.35 a reflexive, symmetric, antisymmetric relation on $\N$ is the equality relation, so there is exactly one.

\sol{5.E.17} False: on $\{0,1,2\}$ the relation with graph $\{(0,1),(1,0),(1,2)\}$ is not symmetric ($1\mathrel R2$ but not $2\mathrel R1$) and not antisymmetric ($0\mathrel R1$, $1\mathrel R0$, $0\ne1$).

\sol{5.E.19} True: $\varnothing\subseteq\N\times\N$, so by Theorem 5.1.14 it is the graph of a unique relation on $\N$, the empty relation $\varnothing_{\N,\N}$.

\sol{5.E.20} True: $\mathrm{Gr}(f)\subseteq X\times X$, so it is the graph of a relation on $X$ by Theorem 5.1.14 (namely $R_f$, Exercise 5.1.11).

\sol{5.E.21} False: for $X=\{0\}$ the empty relation has graph $\varnothing$, but the graph of a function $\{0\}\to\{0\}$ must contain a pair $(0,y)$ (Theorem 3.1.9). (The discrete relation on $\{0,1\}$ is another counterexample: $(0,0)$ and $(0,1)$ both in the graph.)

